\section{Cálculo de Tamanho de Amostra}

O tamanho de amostra necessário para o experimento foi estimado com base na métrica primária (proporção de CPFs com alta taxa de leitura).

Para o cálculo, foi utilizada a calculadora de Evan Miller:

\url{https://www.evanmiller.org/ab-testing/sample-size.html}

Os parâmetros foram definidos a partir da análise exploratória realizada no notebook \url{07_proporcao_cpfs_taxas.ipynb}.

\begin{itemize}
    \item \textbf{Baseline (taxa atual)}: $70\%$, correspondente à proporção observada de CPFs com taxa de leitura $\geq 90\%$ (69.76\% em 268{,}436 CPFs analisados).

    \item \textbf{Minimum Detectable Effect (MDE)}: $5$ pontos percentuais, representando um aumento de $70\%$ para $75\%$.

    \item \textbf{Nível de significância}: $\alpha = 0.05$

    \item \textbf{Poder estatístico}: $1 - \beta = 0.80$
\end{itemize}

Com esses parâmetros, o tamanho de amostra estimado foi de aproximadamente:

\[
n \approx 1{,}335 \text{ CPFs por grupo}
\]
